
\subsubsection{Учет рефракции на мелководье}
\label{sec:wave_refraction}

В предыдущем разделе был рассмотрен механизм шоалинга -- изменения высоты волны при уменьшении глубины, основанный на законе сохранения потока энергии вдоль луча за счет уменьшения глубины в прибрежной зоне. 
Наложив на шоалинг критерий разрушения волны мы фактически определили высоту волн, подходящих к прибрежной точке. 
Однако, этого недостаточно, чтобы вычислить реальное воздействие волн на берег.
Для этого необходимо дополнительно учесть, под каким углом к берегу будут подходить волны перед тем как разбиться об него, передав тем самым свою волновую энергию суше. 

Естественным и наивным решением этой задачи может показаться прямое определение разницы азимутов между нормалью к береговой линии и направлению ветра.
К сожалению, такой подход на практике не является хорошим приближением к реальному механизму волнового движения. 
В реальности, по мере входа на мелководье на волну волну влияет не только вертикальная компонента изменения глубины (приводящая к шоалингу), но и изменения глубин проходящие поперек направления движения волн.
Такие изменения приводят к тому, что участки фронта волны, попадающие на меньшие глубины, замедляются раньше, чем участки над более глубокими местами. 
Аппелируя к жизненному опыту, можно объяснить физическую суть этого явления по аналогии с движущимся автомобилем во время того, когда он цепляет обочину: одно колесо едет по песку (медленно), а другое по асфальту (быстро), из-за чего машина стремится развернуться в сторону более \enquote{тормозящей} среды.
Аналогично и фронт волны начинает разворачиваться, стремясь выровняться по изобатам и принять более перпендикулярное направление движения к береговой линии. 

Этот процесс известен как рефракция волн и представляет собой волновой аналог оптического преломления, происходящего со светом при переходе в среду с отличающейся плотностью.
С точки зрения расчета волновой энергетики, учет рефракции критичен потому, что именно он определяет, под каким углом волновой фронт подходит к берегу и, следовательно, какая доля потока энергии направлена по нормали к берегу, а какая проходит вдоль него. 
Даже если высота волны у берега одна и та же, при разном угле подхода фактическая нагрузка на береговой профиль (и потенциал береговых деформаций) могут отличаться в разы. 

В разрабатываемом методе корректировка направления прихода волны к берегу определялась по закону Снеллиуса.
Современная оптика и гидродинамика интерпретируют закон Снеллиуса как универсальное следствие сохранения фазового фронта и фазовой скорости на границе сред.
Так для света он связывает углы и показатели преломления, а для поверхностных волн на воде -- углы и фазовые скорости, зависящие от глубины ~\cite{LopezRuiz2015_wave_refraction, JooHaJunEffectOfSeaLevel}.

Для двух областей с фазовыми скоростями $C_1$ и $C_2$ и углами $\theta_1$ и $\theta_2$ между направлением распространения волны и нормалью выполняется соотношение:
\begin{equation}
    \frac{\sin\theta_1}{C_1} = \frac{\sin\theta_2}{C_2},
\end{equation}
где $\theta_1$ -- угол волны к нормали изобат в области «глубокой» воды,
$\theta_2$ -- угол в области мелководья у берега.

В мелководном приближении $C_1 = \sqrt{g h_1}$, $C_2 = \sqrt{g h_2}$.
Откуда:
\begin{equation}\label{eq:snellius_low}
    \sin\theta_2 = \sin\theta_1 \cdot \frac{C_2}{C_1}.
\end{equation}

Если $C_2 < C_1$, то $\sin\theta_2 < \sin\theta_1$, а значит $\theta_2 < \theta_1$ и волна повернет ближе к нормали.

Раскрыв формулу~\eqref{eq:snellius_low}, используя введеные ранее обозначения, получим выражение для вычисления угла между нормалью к береговой линии и направлением подхода волны $\theta_\text{out}$ как арксинус величины вычисляемой по формуле:
\begin{equation}\label{eq:snellius_low_upt}
    \sin \theta_\text{out} = \sin \theta_\text{in}\,\sqrt{\dfrac{h_\text{point}}{h_\text{deep}}},
\end{equation}
где $\theta_\text{in}$ задается как угол между азимутом ветра, определяющим направление волны на глубокой воде, и нормалью к берегу.

В случае мелководного приближения $C = \sqrt{g h}$ при $h \to 0$ даст $C \to 0$, что в формуле~\eqref{eq:snellius_low} приведет к тому, что $\theta_2 \to 0$, а волна полностью выпрямляется к нормали.
Однако, так же как и с учетом шоалинга, сама модель не предназначена для описания предельного перехода «на берег», так как в реальности раньше сработает механизм разрушения волны~\eqref{eq:break_wave_eq} и простая формула скорости $\sqrt{g h}$ перестанет отражать реальное поведение процесса.

Поэтому, аналогично формуле~\eqref{eq:green_low_H} в качестве глубины $h_\text{point}$ используется первая от берега положительная глубина, получаемая из батиметрических данных вдоль исследуемого направления.
При этом из-за погрешностей данных батиметрии возникает некоторый \enquote{шум} относительно нуля, что может дать малые положительные значения при интерполяции.
Вычислительно решить эту проблему можно, введя нижнюю границу используемой глубины (например 0.1~м), что даст «минимальную рабочую глубину» для рефракции, ниже которой нет смысла применять мелководную формулу, считая, что волна проходит последний отрезок на той же скорости и в том же направлении, что и при этой небольшой глубине.

Из-за накопления ошибок округления результат вычисления по выражению~\eqref{eq:snellius_low} может получиться выходящим из диапазона значений $[-1, 1]$. 
Для того, чтобы последующий вызов функции арксинуса не выдал комплексное значение, вычисленные по формуле~\eqref{eq:snellius_low_upt} значения жестко обрезаются по интервалу $[-1, 1]$.

Описывая ограничения использованного метода учета рефракции следует отметить, что закон Снеллиуса опирается на линейную теорию волн и лучевое приближение, поэтому его применимость ограничена случаями плавно меняющейся батиметрии, относительно небольших амплитуд и слабых диссипативных эффектов. 
В модели глубина учитывается вдоль одномерного луча, что позволяет описывать поворот волновых лучей к нормали при почти параллельных изобатах, но не учитывает двумерные эффекты дифракции, фокусировки и сложной кривизны береговой линии. 
Задаваемая выражением $C = \sqrt{g h}$ фазовая скорость не учитывает полное дисперсионное соотношение, что делает оценку рефракции приближенной в переходной области глубин. 
Тем не менее, результаты представленные в работе~\cite{LopezRuiz2015_wave_refraction} позволяют сделать вывод о хорошем согласии результатов, получаемых данным методом, с результатами более сложных моделей и использовать этот метод учета рефракции в решении рассматриваемых задач.
